\section{Pile e code}
Le \textbf{pile}, o stack, sono insiemi dinamici dove la posizione di un elemento inserito o rimosso è predeterminata. In particolare, questa struttura
dati implementa la logica \textbf{LIFO}: Last-in, First-out.\par
L'operazione di inserimento è chiamata \textbf{push()} e pone un elemento in cima alla pila; quella di rimozione è invece detta \textbf{pop()} e
rimuove l'elemento inserito più recentemente. Si presentano inoltre i campi \textit{S.top}, l'indice della cima, e \textit{S.size}, la dimensione della
pila. È possibile implementare questa struttura dati su un array di dimensione $S[1:S.top]$.\par
Se la stack è vuota, verificabile con la procedura \textbf{isEmpty()}, allora \textit{S.size = 0} e se si cerca di estrarne un elemento, si avrà un
\textbf{underflow}, mentre se si cerca di inserire un elemento in una pila piena, quindi \textit{S.top = S.size}, si avrà un \textbf{overflow}. Inoltre,
la logica LIFO può essere descritta tramite sequenze di push tranne per i casi in cui si presenta un errore. Infatti, definendo un'operazione
\textbf{top()} per recuperare l'elemento in cima alla pila abbiamo: \[
	Pop(Push(S,x)) = S,\quad Top(Push(S,x)) = x\] \[Top(Empty(S)) = ERR,\quad Pop(Empty(S)) = ERR
\]
\noindent In ogni caso, ognuna delle tre operazioni è eseguita in tempo costante $O(1)$.
\begin{algorithm}[h]
	\caption{isEmpty(S)}
	\begin{algorithmic}[i]
		\Require $S$: La pila
		\If{S.top = 0}
		\State return TRUE
		\EndIf
		\State return FALSE
	\end{algorithmic}
\end{algorithm}
\begin{algorithm}[h]
	\caption{push(S,x)}
	\begin{algorithmic}[i]
		\Require $S$: La pila, $x$: Elemento da inserire
		\If{S.top = S.size}
		\State error "overflow"
		\Else S.top $\gets$ S.top+1
		\State $\text{S[S.top]}\gets x$
		\EndIf
	\end{algorithmic}
\end{algorithm}
\begin{algorithm}[h]
	\caption{pop(S)}
	\begin{algorithmic}[i]
		\Require $S$: La pila
		\If{isEmpty(S)}
		\State error "underflow"
		\Else S.top $\gets$ S.top-1
		\State return S[S.top+1]
		\EndIf
	\end{algorithmic}
\end{algorithm}

\noindent Le \textbf{code}, come le pile, sono un insieme dinamico di elementi con posizioni predeterminate. La differenza principale è
l'implementazione della logica \textbf{FIFO}, First-in, First-out. Ciò significa che gli oggetti sono inseriti nella struttura all'ultima posizione con
\textbf{enqueue()} e prelevati a partire dalla cima con \textbf{dequeue()}. Si possono implementare con un array $Q[1:n]$, il quale avrà gli attributi
$Q.size$ per la dimensione, mentre $Q.head$ e $Q.tail$ sono usati per puntare al primo e all'ultimo elemento rispettivamente.\par
All'inizio, con la coda vuota, i campi di testa e coda sono uguali; provando a usare dequeue() si avrà un errore di underflow, mentre se si presenta una
di queste due istanze: \[Q.head = Q.tail+1 \quad\lor\quad (Q.head=1\land Q.tail=Q.size)\]
\noindent La coda sarà piena, e invocando enqueue() si avrà un errore di overflow. Inoltre, come per le pile, le seguenti operazioni hanno tempo di
esecuzione costante $O(1)$:
\begin{algorithm}[h]
	\caption{enqueue(Q, x)}
	\begin{algorithmic}[i]
		\Require $Q$: La coda, $x$: L'elemento da incodare
		\If{Q.head = Q.tail+1 \textbf{OR} (Q.head = 1 \textbf{AND} Q.tail = Q.size)}
		\State error "overflow"
		\EndIf
		\State Q[Q.tail] $\gets$ x \Comment{x è ora l'elemento in ultima posizione}
		\If{Q.tail = Q.size}
		\State Q.tail $\gets$ 1
		\Else Q.tail $\gets$ Q.tail+1
		\EndIf
	\end{algorithmic}
\end{algorithm}
\begin{algorithm}[h]
	\caption{dequeue(Q)}
	\begin{algorithmic}[i]
		\Require $Q$: La coda
		\If{Q.head = Q.tail}
		\State error "underflow"
		\EndIf
		\State x $\gets$ Q[Q.head] \Comment{Prendi l'elemento in cima}
		\If{Q.head = Q.size}
		\State Q.head $\gets$ 1
		\Else Q.head $\gets$ Q.head+1
		\EndIf
		\State return x
	\end{algorithmic}
\end{algorithm}

%

\section{Liste concatenate}
Una \textbf{lista concatenata} è una struttura dati i cui oggetti sono disposti in un determinato ordine lineare, determinato da un puntatore in ogni
suo elemento. È pratica comune inserire un campo \textit{key} negli oggetti della lista per poterne più facilmente definire l'ordinamento. Esistono più
tipi di liste concatenate: \begin{itemize}
	\item \textbf{Liste singolarmente concatenate}: Gli elementi della struttura sono collegati da un solo puntatore \textit{next}, che indica il nodo
	successivo.
	\item \textbf{Liste doppiamente concatenate}: Gli elementi sono collegati con due puntatori \textit{next, prev}, che indicano il nodo successivo e
	precedente.
	\item \textbf{Liste di ricerca}: Sfruttano il campo chiave per la ricerca efficiente del campo richiesto.
	\item \textbf{Liste ordinate}: Sfruttano il campo chiave per determinare l'ordine dei nodi.
	\item \textbf{Liste non ordinate}: Preservano l'ordine di inserimento degli oggetti.
\end{itemize}

\noindent Contenendo puntatori ai nodi adiacenti, è evidente come si possono definire più istanze di nodo: se il puntatore al precedente è $NIL$,
l'elemento della lista sarà la \textbf{testa}, mentre se è quello al successivo ad esserlo, ci troveremo nell'oggetto di \textbf{coda}. Qualunque altro
elemento sarà intermedio a questi due. Le operazioni attuabili su questa struttura dati sono quelle di ricerca $\Theta(n)$, inserimento $O(1)$ e
rimozione $O(1)$ con chiave, mentre $\Theta(n)$ senza. Si definiscono come segue: \begin{algorithm}[h]
	\caption{list-search(L, k)}
	\begin{algorithmic}
		\Require $L$: Lista concatenata, $k$: Chiave da cercare
		\State $x\gets head[L]$
		\While{$x\neq NIL$ \textbf{AND} $key[x]\neq k$}
			\State $x\gets next[x]$
		\EndWhile
		\State \textbf{return} $x$
	\end{algorithmic}
\end{algorithm}

\begin{algorithm}[h]
	\caption{list-prepend(L, x)}
	\begin{algorithmic}
		\Require $L$: Lista concatenata, $x$: Nodo da inserire
		\State $next[x]\gets head[L]$
		\State $prev[x]\gets NIL$
		\If{$head[L]\neq NIL$}
			\State $prev[head[L]]\gets x$
		\EndIf
		\State $head[L]\gets x$
	\end{algorithmic}
\end{algorithm}

\begin{algorithm}[h]
	\caption{list-insert(x, y)}
	\begin{algorithmic}
		\Require $x$: Nodo da inserire dopo $y$, $y$: Nodo corrente
		\State $next[x]\gets next[y]$
		\State $prev[x]\gets y$
		\If{$next[y] \neq NIL$}
			\State $prev[next[y]]\gets x$
		\EndIf
		\State $next[y]\gets x$
	\end{algorithmic}
\end{algorithm}

\begin{algorithm}[h]
	\caption{list-remove(L, x)}
	\begin{algorithmic}
		\Require $L$: Lista concatenata, $x$: Nodo da rimuovere
		\If{$prev[x]\neq NIL$}
			\State $next[prev[x]]\gets next[x]$
		\Else $\text{ }head[L]\gets next[x]$
		\EndIf
		\If{$next[x]\neq NIL$}
			\State $prev[next[x]]\gets prev[x]$
		\EndIf
	\end{algorithmic}
\end{algorithm}

\noindent In particolare, per quest'ultima operazione, è possibile applicare la \textbf{tecnica delle sentinelle} per semplificare il codice, ignorando
le condizioni in cima e in fondo alla lista. Si tratta dell'implementazione agli estremi di oggetti della lista che rappresentano $NIL$, ma che allo
stesso tempo hanno anche tutti gli attributi dei nodi effettivi, creando quella che è di fatto una lista concatenata circolare con sentinelle.\newpage

%

\section{Alberi}
Le liste concatenate sono un ottima struttura dati per rappresentare relazioni di ordine totale, ma non tutti i casi di studio richiedono che la
proprietà sia estesa all'intero insieme. Introduciamo quindi gli \textbf{alberi}, i quali saranno rappresentati mediante strutture dati concatenate.\par
Gli alberi sono composti di \textbf{nodi}, i quali rappresentano gli oggetti con i relativi campi, e \textbf{archi}, che indicano le connessioni tra i
nodi, implementate con i relativi puntatori in base al tipo di albero scelto. La struttura dell'albero mostra anche un nodo in cima alla rappresentazione,
detto \textbf{radice}, dal quale si diramano i \textbf{figli}, fino a raggiungere la base, i cui nodi sono detti \textbf{foglie}. L'implementazione
generale di questa struttura dati consente a ogni nodo di avere un numero indefinito di figli, ma noi ci concentreremo su tre sottocasi particolari: gli
alberi binari, gli RB-alberi e gli alberi binomiali.

\subsection{Alberi binari di ricerca}
Con \textbf{Albero binario} intendiamo un albero i cui nodi possono avere al più due figli e presentano almeno quattro campi fondamentali: \textit{key},
il valore determinante la relazione di ordinamento, e i tre puntatori al genitore \textit{parent[x]}, figlio destro \textit{right[x]} e figlio sinistro
\textit{left[x]} rispettivamente. La posizione dei nodi deve inoltre rispettare una proprietà dell'ordinamento: i figli sinistri devono avere un valore
chiave minore o uguale del genitore, mentre quelli di destra devono presentarlo maggiore.

\begin{figure}[h]
	\centering
	\begin{tikzpicture}[
		level distance=1.5cm,
		level 1/.style={sibling distance=4cm},
		level 2/.style={sibling distance=1.5cm}
		]			
		\node {$6$}
			child {node {$5$}
			child {node {$2$}}
			child {node {$5$}}
		}
			child {node {$7$}
			child {node {$8$}}
		};
	\end{tikzpicture}
	\caption{Albero binario di ricerca}
\end{figure}

\noindent Per quanto riguarda gli alberi binari \textbf{completi}, ovvero con il numero massimo di nodi foglia ottenibili, le operazioni di ricerca
predecessore, successore, inserimento e rimozione richiedono un tempo $\Theta(log(n))$ nel peggiore degli scenari, mentre nello specifico caso in cui si
crei una catena lineare di $n$ nodi, il tempo peggiora a $\Theta(n)$. Inoltre, l'altezza dell'albero binario attesa è di $O(log(n))$\footnote{Rammento
che, in quanto l'albero si estende esponenzialmente verso il basso creando $n = 2^l$ nodi, il totale dei livelli ottenuti è dato dall'esponente, quindi
$log(n) = l$.}.\par
L'esplorazione dei nodi dell'albero avviene a partire dalla radice e procede muovendosi figlio per figlio; questa operazione è necessaria sia per
ricavare una determinata chiave, sia per l'inserimento o rimozione degli oggetti. In particolare, per queste ultime, si finirà inevitabilmente per
\textbf{sbilanciare} l'albero, poiché si modifica il suo numero totale di nodi; ciò non è un problema perché esistono tecniche per mantenere alberi
sufficientemente bilanciati (con altezza $O(log(n))$) e garantire tempo di esecuzione logaritmico per le operazioni di estrazione e ricerca. Anzitutto,
esistono tre algoritmi principali per l'esplorazione di un albero binario: \begin{itemize}
	\item \textbf{preVisit()}: Visita il nodo corrente ed effettua una visita preliminare dei due nodi figli.
	\item \textbf{inVisit()}: Visita il nodo corrente dopo il suo figlio sinistro e prima di quello destro.
	\item \textbf{postVisit()}: Visita i figli prima del nodo corrente.
\end{itemize}

\begin{minipage}[h]{0.31\linewidth}
	\begin{algorithm}[H]
		\caption{preVisit(x)}
		\begin{algorithmic}[1]
			\Require $x$: Nodo
			\If{$x \neq NIL$}
				\State $visit(x)$
				\State $preVisit(left[x])$
				\State $preVisit(right[x])$
			\EndIf
		\end{algorithmic}
	\end{algorithm}
\end{minipage}
\hspace{0.001\textwidth}
\begin{minipage}[h]{0.3\linewidth}
	\begin{algorithm}[H]
		\caption{inVisit(x)}
		\begin{algorithmic}[1]
			\Require $x$: Nodo
			\If{$x \neq NIL$}
				\State $inVisit(left[x])$
				\State $visit(x)$
				\State $inVisit(right[x])$
			\EndIf
		\end{algorithmic}
	\end{algorithm}
\end{minipage}
\hspace{0.001\textwidth}
\begin{minipage}[h]{0.32\linewidth}
	\begin{algorithm}[H]
		\caption{postVisit(x)}
		\begin{algorithmic}[1]
			\Require $x$: Nodo
			\If{$x \neq NIL$}
				\State $postVisit(left[x])$
				\State $postVisit(right[x])$
				\State $visit(x)$
			\EndIf
		\end{algorithmic}
	\end{algorithm}
\end{minipage}
\vspace{0.01\linewidth}

\noindent Supponiamo ora il seguente albero binario dove sono effettuate le operazioni di preVisit, inVisit e postVisit:
\begin{figure}[h]
	\centering
	\begin{tikzpicture}[
		level distance=1.5cm,
		level 1/.style={sibling distance=4cm},
		level 2/.style={sibling distance=1.5cm}
		]			
		\node {$7$}
			child {node {$3$}
			child {node {$1$}}
		}
			child {node {$10$}
			child {node {$9$}}
			child {node {$12$}}
		};
	\end{tikzpicture}
	\caption{Albero binario di ricerca}
\end{figure}

\noindent I risultati degli algoritmi saranno: \begin{itemize}
	\item \textbf{preVisit()}: $7, 3, 1, 10, 9, 12$.
	\item \textbf{inVisit()}: $1, 3, 7, 9, 10, 12$.
	\item \textbf{postVisit()}: $1, 3, 9, 12, 10, 7$.
\end{itemize}

\noindent Ne risulta che con l'algoritmo \textit{inVisit()} è possibile ordinare gli elementi dell'albero in tempo lineare. Inoltre, la costruzione di
un albero binario richiede una complessità $O(nlog(n))$; ciò vale sia per quelli bilanciati che non. In particolare, per i primi è necessario ordinare
l'array di partenza, ricavarne il mediano per usarlo come radice e inserire poi ricorsivamente il resto dei nodi in base alla proprietà dell'albero.
\begin{algorithm}[h]
	\caption{makeTree(A, p, r)}
	\begin{algorithmic}[1]
		\Require $A$: Array ordinato, $p$: prima metà, $r$: seconda metà
		\State $n\gets r-p+1$ \Comment{Dimensione dell'array}
		\If{$n < 1$} \Comment{Verifica se l'array è vuoto}
			\State \textbf{return} $NIL$
		\EndIf
		\State $m\gets (p+r)/2$ \Comment{Elemento mediano dell'array}
		\State $Node \text{ }x\gets new Node(A[m])$ \Comment{Aggiungi il nodo all'albero}
		\State $left[x]\gets makeTree(A, p, m-1)$
		\State $right[x]\gets makeTree(A, m+1, r)$
		\State \textbf{return} $x$
	\end{algorithmic}
\end{algorithm}

\subsection{RB-alberi}
Gli \textbf{Alberi rosso-neri} o RB-alberi sono strutture dati basate sugli alberi binari di ricerca, il cui scopo è bilanciare l'albero a sufficienza
per mantenere un'altezza logaritmica e quindi garantire l'esecuzione in tempo $O(log(n))$ delle operazioni di base nel caso peggiore.\par
Presentano un campo aggiuntivo apposito per la colorazione di ogni oggetto, il quale può essere rosso o nero. Aggiungendo vincoli in merito alle
modalità di colorazione, è possibile garantire che nessun cammino semplice da radice a foglia sia lungo più del doppio di qualsiasi altro, rendendo
l'albero approssimativamente bilanciato. L'altezza di un RB-albero contenente $n$ nodi è infatti al più $2log(n+1)$, che sta sotto $O(log(n))$.
Questa struttura dati deve inoltre rispettare le cinque proprietà degli RB-alberi: \begin{itemize}
	\item Ogni nodo è rosso oppure nero.
	\item La radice dell'albero è sempre nera.
	\item Ogni foglia (NIL) è sempre nera.
	\item Se un nodo è rosso, i suoi figli sono neri e può essere circondato solo da altri neri.
	\item Per ogni nodo, tutti i cammini semplici dallo stesso alle sue foglie discendenti contengono lo stesso numero di nodi neri.
\end{itemize}

\begin{comment}
	Grazie alla proprietà 3 abbiamo che negli RB alberi avremo al peggio un cammino radice-foglia con alternazioni di nero-rosso fino alla eventuale
	foglia nera, mentre il cammino ottimo avrà esclusivamente nodi neri. Ciò significa che se il meglio è 1, il peggio sarà il suo doppio, 2.
	Manterremo quindi una profondità logaritmica.
	
	Esempio: pag. 300 figura a
	
	RB ha profondità logaritmica e che la ricerca di un elemento specifico, elemento minimo e massimo é logaritmica. Costruzione albero RB non puó
	essere fatta in meno di nlogn.
	'Altezza Nera' BH di un nodo X: Numero di nodi neri incontrati in un cammino da nodo X ad una foglia escludendo X
	I nodi foglia all'estremo con elemento null sono sprecati; consigliato usare il metodo delle sentinelle. Nel caso in cui serva invocare il parent
	di questo nodo nullo, nell'algoritmo invocato bisognerá assegnare al nodo nullo il parent corretto
	
	In generale meglio assumere che la radice sia nera
	Se non lo è la si puó modificare senza alterare nessuna proprietà
	
	Dimostrazione che la profondità è logaritmica.
	Lemma: un RB albero con n nodi interni ha altezza al più 2log(n+1). Un nodo interno è tale se non è una foglia; l'altezza è il cammino peggiore.
	Sottolemma: Il sottoalbero radicato in un nodo x contiene almeno 2^{bh(x)}-1 nodi interni. Quindi prendo un nodo x dell'albero e la considero come 
	radice del sottoalbero.
	
	Dimostrato per induzione sull'altezza del nodo x
	Caso base: altezza 0
	Necessariamente la radice del sottoalbero è una foglia. La sua altezza nera è quindi 0. bh(x)=0. Nodi interni=0. 2^{bh(x)}-1 = 0
	Caso induttivo: altezza h+1
	Abbiamo radice e almeno due sottoalberi; il sottolemma vale per induzione. Tutti i nodi nel sottoalbero di x avranno altezza h.
	Sappiamo poi che per costruzione del RB albero:
	bh(left(x)) >= bh(x)-1; bh(right(x)) >= bh(x)-1
	Interni(x) = interni(left(x)) + interni(right(x)) + 1		// Giustamente il numero di nodi interni comprende gli interi due sottoalberi.
	Interni(x) >= 2^{bh(x)-1}-1 + 2^{bh(x)-1}-1 + 1 = 2^{bh(x)}-1	// Dimostrazione funzionante.
	
	Se l'albero ha profondità h e quindi la radice x ha altezza h, allora sappiamo che l'altezza nera è almeno la metà, quindi bh(x) >= h/2
	Dunque abbiamo detto che la relazione fra il numero di nodi n e le altezze è: n >= 2^{bh(x)}-1 >= 2^{h/2}-1
	Da questo possiamo ricavare: n+1 >= 2^{h/2} \implies log(n+1) >= h/2 \implies h = 2log(n+1). La dimostrazione è vera.
	
	Quindi l'altezza di un rb albero con n nodi è al più il doppio di log(n). La profondità è logaritmica.
	
	--- Come inserire un nodo mantenendo profondità logaritmica?
	Partiamo dal presupposto che si segue la logica di inserimento di un qualunque albero binario di ricerca. Quindi si sposterà a destra o sinistra in
	base alle regole di definizione.
	Di che colore facciamo il nodo? Se lo facciamo rosso e siamo una foglia, viola la definizione, mentre se lo facciamo nero sminchia la profondità
	logaritmica.
	Quindi lo facciamo rosso a prescindere lmao
	Abbiamo un rb albero con l'anomalia sul nodo appena inserito. Bisogna eliminarla.
	Facciamo delle trasformazioni locali portando l'anomalia verso la radice. Oppure la eliminiamo prima, dipende quale vien prima. Se questo problema
	si presenta alla radice, non è un problema. È valido.
	
	Domande:
	- Possiamo assumere che la radice sia sempre nera? Sì, perché se la radice fosse rossa potremmo colorarla di nero senza eludere la definizione.
	- Concetto di rotazione: nell'operarci cambieremo e ruoteremo i nodi dell'albero.
	Ruotare nel senso scambiare il babbo col figlio. Si può fare a dx e a sx
	
	--- Algoritmo di ribilanciamento; rotazione
	
	Pseudo codice
	P sta per parent
	Inoltre assumi che la radice sia Nera
	
	L'ultimo Else é di if(color(P[x])== red)
	
	Il caso 1 propaga verso l'alto
	
	Mentre il caso 3 é concatenato al 2 e termina il processo
\end{comment}

\begin{comment}
	Algoritmi - 01/12/25
	
	Inserimento estrazione e ricerca in tempo logaritmico se si usa un RB albero
	Inserimento, con eventuale anomalia, la quale viene eliminata o propagata verso l'alto
	Gli algoritmi per le rotazioni hanno complessità costante. L'unico caso con complessità maggiore è quello dove si spinge l'anomalia verso la
	radice, ma è di massimo logn
	
	E' possibile anche unire due RB-alberi che contiene gli elementi di entrambia? Sì, ma in minimo tempo lineare, prima con una buildtree e poi un
	merge.
	
	Ex. Dato un nodo x, trovare il suo successore (quello che nell'ordinamento occuperebbe la posizione successiva)
	Se esiste il figlio destro, prendo il minimo del figlio destro, mentre in loro assenza, si guarda agli antenati destri.
	
	next(x) {
		if (right(x) != NIL && x != NIL)
		ret min(right(x))
		else {
			while (x == right(p(x)) && p(x) != NIL) {
				x <- p(x)
			}
			if (x != root) return p(x)
			else return NIL
		}
	}		// Complessità nlogn, non supero la profondità dell'albero.
	Le next hanno complessità lineare, tranne nel caso pessimo che è nlogn, dove deve scorrere l'intero albero.
	
	Giustamente puoi farci anche la select. Il minimo e il massimo son banali, alle estremità sinistra e destra.
	Tuttavia considera il mediano ed un albero sbilanciato. Per un tempo logaritmico bisogna seguire un singolo cammino. Ciò si può ottenere
	aggiungendo un campo "rank" ad ogni nodo.
	Questo consentirebbe di capire la direzione istantaneamente, usando un singolo cammino.
	However; se vado ad inserire un nuovo minimo, toccherà cambiare il rank di tutto l'albero, e ciò ha una complessità lineare. non sufficiente.
	
	Al posto di rank diciamo invece un campo size, che dirà quanti nodi ci sono nell'albero radicato in x.
	Vediamo come size possa essere mantenibile per inserimento ed estrazione:
	- insert: il size aumenta esclusivamente nei nodi toccati dal percorso.
	- extract: Stessa cosa; rimuovendo il nodo in un cammino, verrà modificato solo il size dei nodi del percorso.
	- Occhio che nelle rotazioni il size cambia, sebbene la loro somma rimanga uguale.
	Togo, il campo size quindi non va ad influire sulla complessità delle operazioni elementari, facendole rimanere tutte a complessità logn.
	Also, capiamo: size[x] = size[left(x)] + size[right[x]] + 1
	
	// select logaritmica con il campo size
	select (x, i) {
		if (i >= 1 && i <= max) {
			r <- size[left(x)]+1
			
			if (r == i) ret x;
			else {
				if (i > r) ret select(left(x), i);
				else if (i < r) ret select(right(x), i-r);		// i-r perché nel subtree destro, gli elementi di r non ci sono più.
			}
		}
	}
	
	Se volessimo una posizione specifica in senso di rank, tuttavia, questa cosa funzionerebbe solo per la radice, quindi bisogna trovare un modo per
	generalizzarlo
	
	rank (x) {
		r <- size[left(x)]+1;
		y <- z;
		
		while (y != root) {
			if (y == right(P(y)) {
				r += size[left(P(y))]+1;
			}
			y <- p(y);
		}
		ret r;
	}
	
	Dovessi aggiungere un campo ai nodi dell'rb albero, è necessario dover dimostrare che mantenga la complessità logaritmica per le operazioni.
	What if esistesse un modo per determinarlo senza volersi ammazzare? Lmao esistone il teorema, basta attenersi alle sue regole.
	
	Sia F[x] un campo aggiuntivo. Se esiste una funzione f(x) tale che:
	F(x) = f(key(x), F(left(x)), F(right(x)), key(left(x)), key(right(x)))
	allora F è mantenibile. La potenza del teorema cresce in base a quante operazioni ha.
\end{comment}


\subsection{B-alberi binomiali}

%

\section{Estensione di strutture dati}

\begin{comment}
	02-12-25 - Algoritmi
	
	--- Estensione di strutture dati
	Prendere per esempio un RB albero ed aggiungerci campi comodi per poter eseguire più facilmente gli ordinamenti. In questo caso specifico la cosa
	più importante è non andare
	a considerare il nodo parent negli argomenti di F(x).
	I campi, infatti, cambiano esclusivamente se viene inserito/rimosso un nodo sullo specifico cammino; gli altri rimangono intonsi. E si fa pure in
	logn sta cosa.
	
	Per le operazioni di risistemazione dell'albero:
	- i colori non cambiano, non è un argomento di F(x).
	- Nelle rotazioni, i nodi ruotati x,y possono avere i loro campi cambiati; quindi costa logn, perché devo propagare verso l'alto il campo.
	
	Dove quest'ultima sfaccettatura è un problema, l'algoritmo rimane ugualmente logaritmico nonostante la rotazione non sia più di costo costante.
	Più precisamente, di base:
	- inserimento: logn + 2rot, nel caso esteso preso in esame 2rot=2logn
	- estrazione: logn + 3rot, " 3rot=3logn
	Ma complessivamente rimaniamo sempre su logn.
	
	// Esempio: agenda con appuntamenti, caratterizzati da un intervallo di tempo, l'obiettivo è ins/ext e trovare conflitti, tutto in tempo logaritmico
	// Serve un modo per decidere istantaneamente se andare a sx o dx. Ordiniamoli in base al tempo e aggiungiamo il massimo tempo di fine visita come
	campo aggiuntivo.
	
	search (x, I)		// I intervallo, tempo di inizio e fine dell'appuntamento. x è la radice.
	while (x != NIL && I \cap int(x) == \void) {			// int è intervallo di x.
		if (left(x) != NIL && max[left(x)] > low[I])		// max è massimo estremo destro del sottoalbero sx, low è l'estremo inferiore di I.
		x <- left[x]
		else 
		x <- right[x]
	}
	ret x;
	
	Supponiamo di avere un determinato intervallo in root. Sappiamo che se non troviamo nulla a sinistra, allora sicuramente non c'è nulla nemmeno a
	destra.
	Se un subintervallo left ha massimo nel mezzo dell'intervallo root, sappiamo che a destra esisterà un altro intervallo che sarà anche dentro questo
	max. Dunque c'è un conflitto.
	Se non ci sono conflitti a sinistra, allora non ce ne saranno neanche a dx.
	
	// Esempio 2: Abbiamo lanciato una missione su marte, dove vengono raccolti vari campioni di suolo. Sono catalogati in un database, ordinato in base ad una chiave, la frequenza
	// spettrale dominante. Il misuratore ha una tolleranza \epsilon; rendendo la frequenza [f+\epsilon, f-\epsilon]. Come memorizzare i dati sicché il
	db restituisca, se esiste un
	// campione corrispondente a quello letto?
	// Può essere che ci siano più frequenze candidate, bisogna trovarne la frequenza più vicina.
	
	search (x, f)		// f è la frequenza.
	if (x == NIL) ret NIL;
	if (f < x.freq) y <- search (left(x), f);
	else y <- search (right(x), f);
	
	if (y == NIL || mod(f-x.freq) < mod(f-y.freq)) return x;		// mod è modulus, il valore assoluto.
	else return y;
	
	bestCompatible(x, f, \epsilon)
	y = search (x, f)
	
	if (y != NIL && mod(y.freq - f) < \epsilon*2)
	ret y
	else
	ret NIL
\end{comment}

%

\section{Heap binomiali}

\begin{comment}
	15-12-25 Algoritmi
	
	Voglio una struttura dati dove si può fare inserimento, estrazione, ricerca e unione in tempo logaritmico.
	Ha stile heap. Si chiama Heap binomiale, ed è un insieme di alberi binomiali, definiti per induzione, ognuno caratterizzato da una data dimensione.
	
	-- Albero binomiale
	B_0 = un singolo nodo
	B_{i+1} = faccio diventare figlio della radice del primo albero l'intero secondo albero.
	
	Hanno alcune proprietà:
	- In un albero binomiale di dimensione k ci sono 2^k nodi, quindi il numero di nodi dell'albero cresce esponenzialmente con la dimensione.
	2^i+2^i = 2*2^i = 2^{i+1}
	
	- L'altezza ha valore k, quindi esattamente pari alla dimensione.
	Quasi ovvia la dimostrazione, basta prendere il percorso più profondo.
	
	- Alla profondità i ci sono (k\\i) nodi.
	E' un triangolo di tartaglia stampato come un array che aumenta di 1dim ogni iterazione:
	1
	11
	121
	1331
	14641 etc.
	k=riga, i=colonna
	
	k=0 + i=0 \implies nodi=1 ok
	k=1 + i=1 \implies nodi=2 ok
	k=k+1 + i=i \implies nodi nell'albero binomiale figlio  traslati di i-1 rispetto alla i dell'albero genitore.
	Quindi (k+1\\i) = (k\\i-1) + (k\\i), grazie alla proprietà triangolo di tartaglia.
	
	- I figli della radice da sx a dx sono B_{k-1} ... B_0. I figli di un albero binomiale sono tutti gli alberi di dimensioni minori.
	
	Un heap binomiale è una lista di alberi binomiali tali che:
	- per ogni dimensione esiste al più un albero.
	- I nodi contengono chiavi che soddisfano le proprietà di un heap.
	
	L'elemento minimo è in una delle radici degli alberi. La complessità della ricerca del minimo è lineare nel numero degli alberi.
	
	Esempio:
	-> 5 -> 9 -> 17
	|	|  |
	11	18 18
	|
	21
	
	-> 7 -> 4 -> 27 -> 25
	|	 /\	   /\
	10	 B_4   B_5				// B_4 e B_5 sono due alberi binomiali di dimensione rispettiva 4 e 5.
	
	Perché non ci possono essere più di logn radici distinte? Qualunque grado dell'albero io prenda deve avere 2^k \leq n. Dunque il grado massimo che
	si può prendere è k \leq log_2n
	E no, non posso andare oltre perché posso inserire al più l'albero binomiale della dimensione più tutte le dimensioni precedenti.
	Dato n converti il suo valore in binario e otterrai facilmente le posizioni degli alberi binari.
	La ricerca del minimo sarà quindi in realtà logaritmica, perché scorri ogni radice.
	
	--- Unione di due heap binomiali
	Somma di due numeri binari; facciamo il merge delle due liste:
	-> 5 -> 7 -> 9 -> 4 -> 17 -> 27 -> 25
	|	  |	   /\    /\    /\
	11	  10   B_2   B_4   B_5
	
	-> 5 -> 9 -> 4 -> 17 -> 27 -> 25		// Fuso 5, 7 in un albero di grado superiore
	|	|	  |	   /\    /\    /\
	7	11	  10   B_2   B_4   B_5
	
	-> 5 -> 4 -> 17 -> 27 -> 25				// Fuso 4, 9 in un albero di grado superiore
	|   / \	 /\    /\    /\
	7  9  10	 B_2   B_4   B_5
	|
	11
	
	-> 5 ->  4 	-> 27 -> 25					// Fuso 4, 17 in un albero di grado superiore. Tutte le dimensioni diverse. Finito.
	|   / /\	   /\    /\
	7  17 9 10  B_4   B_5
	/\  |
	B_2 11
	
	000111 + 110011 = 111010
	
	L'algoritmo ha un pointer in un nodo, il suo successore ed il successore del successore.
	Va a guardare i gradi dei nodi x, del successore e del successore del successore. Se nota gradi uguali, fa una fusione fra i due, altrimenti passa
	ai prossimi.
	
	La complessità dell'unione è logn.
	Per calcolare la dimensione puoi estendere la struttura dandogli un campo size, così la complessità diventa costante.
	
	-- Inserimento
	Prendi l'elemento, vedilo come un heap binomiale con un singolo nodo. Fai l'unione seguendo le regole appena viste.
	
	-- Estrazione
	Per la proprietà 4, se rimuovi la radice di un heap binomiale, rimane quello della sua dimensione-1. Fai una unione e la struttura è nuovamente
	corretta.
	
	-- ReduceKey; voglio un nodo e cambiargli la chiave con un nodo più piccolo
	// Sta roba ha senso farla perché la useremo per eliminare ogni elemento arbitrario di un albero.
	ReduceKey(x, y)
	
	// Riduce il valore di un dato nodo, per poi fare heapify per rimuovere l'anomalia. Sta cosa ha complessità logn.
	// Arriverai che questo nodo sarà la nuova radice con una chiamata ricorsiva a reducekey, per poi rimuoverlo facilmente.
\end{comment}

%

\section{Insiemi disgiunti}

%

\section{Tabelle hash}